% ===========================================================================
% 06-Conclusions.tex. Presentation only; numbers verbatim from v1.
% Terminology, load-bearing: held-out treated WELL, trained CONDITION.
% "treated well" is the glossary's form (\glentry, 00-HowToRead:267) and wins
% over "treatment well" everywhere; this file used the loser three times.
% Never "pairing"; the thesis declines that reading.
% ===========================================================================

\chapter{Prompt sensitivity is established; biological transfer is not}
\label{sec:conclusions}
\framedecl{ER}


\begin{question}
Does a cell-sentence language model fine-tuned on this atlas learn what a
particular drug does to a particular cell? The instrument was a chain built
in order: a metric calibrated before it was trusted, controls a drug-blind
predictor cannot pass, and a bar that needs no model at all. The bar was not
cleared.
\end{question}

Prompt sensitivity to trained drug names is established; biological
transfer is \notestablished, and nothing built here beats the training-only
per-drug lookup on common support.
\noindent
The win condition, fixed in advance in \cref{sec:introduction}, is that a
conditional model beat the training-only per-drug lookup on common support. Only
then has its access to the untreated cell supplied information a dictionary of
per-drug averages lacks. A split-half reference cannot substitute, because it
cannot separate a drug's characteristic response, which transfers, from this
context's departure from it.

The expression-target checkpoint makes no useful condition-specific predictions.
Residual encoding reveals real sensitivity to trained drug names, but the
checkpoint recovers less than a simple lookup and establishes no transfer to new
drug--context combinations or unseen drugs.

That verdict ends a chain, and no isolated point estimate may leapfrog it.

\begin{scopelimit}[carried out of the thesis; see \cref{sec:limitations}]
One $1$B C2S-Scale Pythia backbone; five canonical fine-tuning arms cold-started
from it, at most one epoch at one learning rate (the consensus arm stopped at
$61.5\%$); no larger checkpoint. A separate ten-epoch residual run tests ordinary
undertraining only through held-out cross-entropy; it is not a sixth behavioural
arm, and no long-run checkpoint was evaluated with $\NIR$. These are
checkpoint-level results, and a limit for conditional generative models in
general is \notestablished by them.
\end{scopelimit}

% ---------------------------------------------------------------------------
\section{Only one of five metrics survived its own calibration}

A negative result is worth no more than the ruler that produced it, and the
field's default ruler failed a test it should have passed. A zero-information
predictor placing every gene at the middle rank scores
\ci{\DEdr=0.999913}{0.999908}{0.999919}, above both the two-cell reference and
the model. Dynamic Range Fraction therefore grades a metric by whether it rewards
a real replicate more than that arm.\seealso{\cref{sec:res-metrics} builds the
audit and reports all five arms}

Only $\NIR$ of the five is positive, returning \ci{+0.635}{+0.604}{+0.663} across
plates over $25$ cell lines and \ci{+0.545}{+0.516}{+0.569} within plate over
$27$, each at Holm $p=0.002$. Three of the other four invert; the weighted $R^2$
variant lies nearest zero with its sign tracking cell count, so no directional
claim is made for it. A residual-frame calibration returns
\ci{+0.704}{+0.667}{+0.736} over $42$ cell lines against a zero-drug-information
control. All are percentile cluster bootstraps without small-sample correction,
clustered on cell line after an earlier artifact clustered on drug
rows;\corrmark{corrected in \cref{sec:res-metrics}} plate stays crossed with cell
line, so they are robust on that axis alone.

What travels is the procedure. A perturbation metric must first show positive
dynamic range against a real replicate and an uninformative predictor; the lookup
is a second such bar, the removal gate a third.

\paragraph{Where this leaves the calibration dispute} Miller et
al.~\cite{miller} argue that miscalibrated metrics can manufacture apparent
failures of deep perturbation models (\cref{sec:lit-dispute}). This thesis
confirms the point, applies it, and still reaches a different verdict. Their
rank-based $\NIR$ survives the audit and the checkpoint is at chance on it, which
bounds their conclusion's extension from genetic perturbation to these chemical
predictions without overturning it.

% ---------------------------------------------------------------------------
\section{In the expression frame the checkpoint is drug-blind}

On the calibrated metric the checkpoint fails useful condition-specific
prediction. Within plate on unseen drugs it scores $0.498$ against $0.504$ for a
control-copy that reproduces the untreated cell and $0.576$ for a within-well
split-half precision reference.

\paragraph{The null is not a failure to measure} A null can always be blamed on
the wells, so the aggregate was re-scored where the wells are demonstrably
adequate. On the $204$ conditions whose own reference clears the pre-set $0.80$
identifiability bar and averages $0.938$, the model reaches $0.768$ and
control-copy $0.766$, a paired difference of \ci{+0.002}{-0.026}{+0.028}.
Split-half retrieval does not prove a drug's response is predictable from an
untreated cell, but at the tested design and cell budget it rules out a failure
to measure distinct treated populations.

\paragraph{Three instruments that fail differently agree} Changing the drug and
mechanism while holding the cell fixed moves the prediction by
\ci{+0.014}{-0.016}{+0.042} on the identifiable subset. Output invariance, which
asks whether changing the drug token changes the generated text at all and not
whether it changes it correctly, returns \ci{0.000}{-0.019}{+0.016} in top-$100$
Kendall $\tau$, bounded at $\leq 0.016$ against $\tau = 0.198$ between two
generations from a byte-identical prompt; the Jaccard version is tighter,
\ci{+0.000}{-0.008}{+0.008}. Forced choice, grading a prediction against its own
condition's truth and against another condition's, returns an own-truth win rate
of $0.48$ against a chance $0.50$ fixed by construction, where the precision
reference reaches $0.67$--$0.83$. Their weaknesses do not overlap, which is why
the agreement carries weight. No drug effect is \shownzero by any of them, and
any effect they miss is too small to support the intended prediction task.

% ---------------------------------------------------------------------------
% Third question of sec:intro-map, restated at the altitude of the other three.
% NB not "The drug is not missing from the computation" and not "Identity is
% read, at very low gain": those are verbatim the beat titles of
% sec:methods-probe and sec:res-mechanism, and duplicate ToC entries read as an
% error. The data branch stays in the drug-blindness section above, where the
% body also puts it (04b-representation.tex:44 defers to sec:res-blind for it);
% the lede below carries it across so all three branches are answered here.
\section{The drug is not absent; it is barely consulted}

The third question of \cref{sec:intro-map} asks where that null comes from: the
drug absent from the data, absent from the model's internal representation, or
present in the representation and not consulted when the model generates. The
re-scoring above rules out the first at the tested design and cell budget.
Behaviour cannot separate the other two, since both predict the same scores, so
the weights were read instead.

\paragraph{The label is not lost before generation} After cell line, plate and
dose are projected out, drug identity becomes more linearly decodable with depth,
reaching $0.88$ at layer $16$ against an $8\%$ chance floor. A later probe suite
asked whether that signal is read; its activation-swap interpretations did not
survive their controls.\corrmark{the swap arm and its failed controls are in
\cref{sec:res-mechanism}} What survives supports only low-gain reading of a
decodable drug label, in one training arm and not another, surviving pooled
multiplicity correction at one displacement rung only and on one of two stored
bootstrap draws; no conclusion rests on it alone.

% ---------------------------------------------------------------------------
\section{Re-encoding produces the first trained-drug effect}

If the drug is in the data and in the activations, the next suspect is the
training target. Standard cell-sentence targets for two drugs in the same context
share $82.7\%$ of their top-$200$ genes, a property of the targets and not a
token-level attribution of the gradient. Residual encoding raises the differing
fraction to $59\%$ at cell-line scope and $60\%$ at the training plate scope.

Against the geometry-defined \texttt{orth} partner, trained conditions score
$+0.0898$, with intervals excluding zero clustered on cell line and well,
\ci{+0.0898}{+0.0382}{+0.1415}, and on drug and well,
\ci{+0.0898}{+0.0291}{+0.1506}. That is the project's first measurable
trained-drug effect.

\paragraph{The comparator is part of the measurement} A prompt replacement is not
automatically a null. The scrambled arm itself scores $0.550$, $0.501$ and
$0.423$ as the partner moves from aligned through orthogonal to anti-aligned, so
a dissimilar lie inflates the apparent benefit of the correct name. The
\texttt{orth} stratum, with the comparator's own score reported, avoids that
inflation without manufacturing an independent null.

\paragraph{The gain identifies a package, not a mechanism} Residual encoding
changes aggregation, control referencing, generic subtraction, sign coding and
sequence length together, and separating them needs a factorial comparison that
was not run. Nor is it uniquely responsible. At the matched \num{1400}-token
budget the optimal-transport arm also clears zero against its orthogonal
comparator, \ci{+0.0292}{+0.0144}{+0.0441}, where its $600$-token estimate
spanned zero. A generation cap can manufacture a null, so any claim about the
residual advantage must retain the budget comparison.

\paragraph{Ordinary undertraining was examined separately} Validation
cross-entropy on the ten-epoch residual diagnostic is best after epoch 1
($1.9199$) and ends at $2.2568$ while training loss keeps falling. The schedule
was stretched, not continued, and no checkpoint from it was scored with $\NIR$;
extending this optimisation does not improve its held-out objective, and no claim
is made that every schedule must fail.

% ---------------------------------------------------------------------------
% NB not "Responding to the drug is not predicting it": verbatim the beat title
% of sec:res-generalisation, and two identical ToC entries read as an error. The
% title below names the effect and the missing transfer instead, so the
% conclusion states the finding rather than echoing the results heading.
\section{The trained-drug effect establishes no transfer}

One question is whether the prediction moves when the drug named in the prompt
changes. The other is whether the prediction is right, moving \emph{toward} the
response of that drug. The trained-condition effect establishes the first for
drugs represented in fine-tuning, and it does not establish biological transfer.

Conditions in held-out treated wells score $+0.0332$, excluding zero clustered
on cell line and well, \ci{+0.0332}{+0.0069}{+0.0594}, but not on drug and well,
\ci{+0.0332}{-0.0031}{+0.0695}. Drug is the axis a transferred residual must
generalise over. That split is largely recombination, with $30$ of its $34$ drugs
and $39$ of its $42$ cell lines in training and $82\%$ of its conditions
recombining a training drug with a training cell line. On genuinely unseen drugs
the gap is \ci{-0.0315}{-0.0836}{+0.0206}, both axes spanning zero. The splits
differ in reliability selection, so their point estimates are not a clean decay
with novelty.

\paragraph{The partner slope measures the trained drugs} It is positive on every
split, but fitted only over scrambled partner prompts, most of which are
themselves trained drugs on the unseen-drug split, and the true target drug is
not a point on that line. Held-out predictions sit near chance when ranked
against other drugs in the same cell line. The map is coarse and
trained-drug-to-residual, not a transferable response model.

% ---------------------------------------------------------------------------
\section{Nothing built here beat the lookup}

A per-drug lookup needs no checkpoint, no prompt and no cell
sentence.\seealso{\cref{sec:res-lookup} scores the full ladder through one
pipeline} It predicts a drug's residual as its mean residual in other cell lines,
fitted from training conditions only.

On the \num{1192} conditions where both are defined it scores $0.913$, the model
$0.549$, the within-well reference $0.857$. It reaches $0.980$ on trained
conditions, $0.890$ on held-out ones, $0.822$ restricted to one other cell line.
Its excess over the split-half reference reflects a many-line average against a
noisy half-sample on a metric compressed near its upper limit, not performance
beyond biological headroom. Nothing conditional built here beats direct retrieval
of the drug's mean residual.

\paragraph{A bar, not a full account} That average is not a full account of the
drug's effect. A disattenuated cross-context transfer coefficient is
\ci{\Tcoef=0.553}{0.514}{0.593}; a fully shared simulated null returns $1.001$
and a structure-matched negative control \ci{-0.008}{-0.036}{+0.020}. About
$45\%$ of the similarity scale on which $\Tcoef$ is measured is not shared across
compared contexts. The lookup's $0.913$ therefore does not mean drug response is
nearly constant; it means $\NIR$ recognises the shared per-drug component and
compresses distinctions near its upper end.

\begin{scopelimit}[carried out of the thesis; see \cref{sec:limitations}]
The unshared part is explicitly \emph{not} a drug-by-cell-line interaction or a
variance share. The estimator mixes cell line, dose, well and target estimation,
was run with \texttt{same\_dose\_only} and \texttt{loo\_generic} off, and still
leaves roughly $0.3$ unshared when dose varies within one cell line. Reading
$1-\Tcoef$ as the learnable opportunity is equally wrong. It is an upper bound on
what the constant leaves unshared under this estimator, not evidence a model can
predict that remainder.
\end{scopelimit}

Two tests looked for learnable structure in that remainder and found none
detectable. Within-line consistency is \ci{-0.0149}{-0.0251}{-0.0050} against
\ci{+0.0003}{-0.0071}{+0.0081} across lines, and a mechanism-matched context test
is null. Neither makes the remainder irreducible. That second quantity has low
split-half reliability and mechanism is heavily confounded with plate, so
resolving it needs more cells per condition or a design crossing mechanism with
plate.

\paragraph{DrEval's lesson, with a vector in place of a scalar} DrEval finds deep
drug-sensitivity models often add little beyond mean drug and cell-line
effects~\cite{dreval}; here the predicted object is a $946$-gene vector, not
scalar potency, and a per-drug mean again dominates. Simple baselines are not
thereby biologically complete. A conditional architecture earns its complexity
only by beating the retrieval component already available from training.

% ---------------------------------------------------------------------------
\section{For an unseen drug, one lead is live on the cell-line axis and nothing
is closed}

No lookup exists for a drug absent from training, so the final gate asks which
shared property could stand in for one, testing three channels with partners
restricted to training conditions.

The gate reports in four states, and the distance between them is the point. A
channel is \textsc{live} when its interval clears the pre-set $0.03$ relevance
margin on the clustering axis the claim rides on, \textsc{closed} when the
interval lies wholly below that margin, and \textsc{inconclusive} otherwise. It
is \textsc{untestable} when the matched control could not be built at all: ``we
could not build the control'' must never be written as
``closed''.\seealso{\cref{sec:res-scope} prints the gate's verdicts beside each
channel's coverage}

\paragraph{The verdicts} Protein target clears the margin under all three
constructions clustered on cell line and well; clustered on drug it clears only
against the weaker count-matched null and falls short against both
estimand-correct nulls. Mechanism excludes zero on the cell-line axis but falls
just below the margin and spans zero on the drug axis. Chemistry is narrower and
unresolved, excluding zero against the plate-matched null on the cell-line axis
but spanning zero in the different-plate construction and under both nulls on the
drug axis.

Cell-line clustering asks whether a channel effect is stable across represented
cellular contexts; drug clustering asks whether it generalises across the
drugs an unseen-drug claim rests on, so target can be live under the first
and inconclusive under the second without contradiction. The three nulls answer
different questions and the arms differ in support and annotation, so their order
is not a biological ranking.

The strongest lead rests on only $18$ target-annotated drugs, two held out.
Protein-target annotation carries transcriptional signal on that mixed-support
arm, but a Leave-Drugs-Out result is \notestablished by it, and mechanism and
chemistry are inconclusive rather than closed. The arm's estimated $80\%$-power
detection threshold is about $0.079$ in $\NIR$, so effects materially smaller
than eight points are not well powered. The next experiment needs enough held-out
drugs, preferably held-out mechanism families, for drug-clustered inference.

\paragraph{Two repairs that must be designed, not assumed} One sentinel token
cannot be expected to map to a transcriptional residual, so a target or mechanism
description is the obvious bridge from an unseen name. It must be read against
the backbone's prior knowledge: holding a drug out of Tahoe fine-tuning does not
erase medicinal-chemistry text from language pretraining. That arm needs matched
description quality, missingness controls, a frozen knowledge probe and enough
held-out drugs for drug-clustered intervals. A larger backbone is a
replication, not a scaling law, unless the family is held fixed, since Pythia-1b
to Gemma-2 changes tokenizer, architecture, corpus and scale together.

% ---------------------------------------------------------------------------
\section{The controls are most of the result}

In a literature where uninformative baselines can score well, a negative result
is worth what its controls are worth. Each surviving claim rests on a control
added after a specific failure. Plate-matched comparison sets followed a
zero-information baseline that improved when plate was free to vary; a
swap-distance sweep followed a weak similar-drug scramble. A removal gate
preceded any ablation null and matched-norm injection any swap interpretation.
Intervals were
clustered on the units the estimand requires, cell line, drug and treated well,
never on cells, and a pre-registered interpretation was allowed to fire its
inconclusive branch.

Claims that failed those controls were withdrawn, and the withdrawals recorded,
not quietly dropped. That record is why the remaining claims can be read at the
strength stated here, and no stronger.

\begin{claim}
On the \num{1192} conditions where both are defined, the training-only per-drug
lookup scores $0.913$ against the model's $0.549$, with the within-well
reference at $0.857$; restricted to a single other cell line it still scores
$0.822$. Residual re-encoding does produce a real trained-condition effect
against the \texttt{orth} partner, with intervals excluding zero on both
clustering axes, and it does not establish transfer: conditions in held-out
treated wells give an interval that excludes zero on the cell-line axis and
spans zero on the drug axis, which is the axis a transferred residual must
generalise over, and on genuinely unseen drugs the gap is
\ci{-0.0315}{-0.0836}{+0.0206}
with both axes spanning zero. No drug effect is \shownzero by any instrument
here; the three that bound it fail differently and agree. Cross-context biology
leaves a real remainder, \ci{\Tcoef=0.553}{0.514}{0.593}, but $1-\Tcoef$ is an
upper bound on what a per-drug constant leaves unshared under this estimator,
not a learnable opportunity. Of the three channels open to a drug with no
lookup, protein target is live on the cell-line axis and inconclusive on the
drug axis on $18$ annotated drugs of which two are held out, mechanism and
chemistry are inconclusive, and none of the three is closed.
\end{claim}
